% Generated draft: requires mathematical and source review.
\documentclass[11pt]{article}
\usepackage[margin=25mm]{geometry}
\usepackage{fontspec}
\setmainfont{Latin Modern Roman}
\usepackage{amsmath,amssymb,mathtools,unicode-math}
\setmathfont{Latin Modern Math}
\usepackage{enumitem,needspace,xurl,hyperref}
\setlength{\parindent}{0pt}
\setlength{\parskip}{5pt}
\setlength{\emergencystretch}{3em}
\newcommand{\N}{\mathbb{N}}
\newcommand{\Z}{\mathbb{Z}}
\newcommand{\Q}{\mathbb{Q}}
\newcommand{\R}{\mathbb{R}}
\newcommand{\C}{\mathbb{C}}
\newcommand{\F}{\mathbb{F}}
\newcommand{\A}{\mathbb{A}}
\newcommand{\PP}{\mathbb{P}}
\newcommand{\E}{\mathbb{E}}
\newcommand{\eps}{\varepsilon}
\newcommand{\dd}{\,\mathrm{d}}
\begin{document}
\section*{UnsolvedMath ID: 11000136}
\textbf{Source problem code:} catalogue entry
\subsection{Definitions and mathematical statement}
Retain \(\operatorname{Arc}(M)\) as the arc complex of a surface in the source's convention. Its vertices, weights, quotient by mapping classes, boundary markings, and admissible surfaces are not defined in the excerpt. In particular, an ordinary unquotiented simplicial arc complex must not be silently substituted for the stratified \(\operatorname{Arc}\) space intended here.

\emph{Problem 6 (R. Penner's chapter).} Calculate the topological type of the arc complexes: PL-homeomorphism type, homotopy type, homology, and related invariants. The first nontrivial cases requested are the PL manifolds \(\operatorname{Arc}(M)\) for four ``type one'' surfaces, which are not identified in the supplied record.

The source proposes this recursive finite-type description of stratified arc-complex topology:
\begin{itemize}
\item A PL sphere is a type zero space.
\item A type one space is one of a list of four specified closed, connected, simply connected, nonspherical manifolds: the arc complexes of the four type-one surfaces, with respective dimensions \(5,7,7,9\). This is a list of four spaces, not every manifold of those dimensions. Their explicit topological types are part of the question.
\item For \(n>1\), a type \(n\) space is a finite polyhedron, considered up to PL isomorphism, whose vertex links in every compatible triangulation are PL isomorphic to iterated suspensions of joins of at most two spaces of types strictly less than \(n\). The admissible zero-factor and zero-suspension conventions are not specified.
\end{itemize}
The accompanying conjecture is that each arc complex in the intended surface class has some finite type:
\[
 \forall F\ \exists n<\infty,\qquad\operatorname{Arc}(F)\text{ has type }n.
\]
The level may depend on \(F\); no uniform type bound is supplied. The source says Theorem 5 establishes the required form for many simplex links, and says a particular collapsing argument in the ``calculus of mapping cylinders'' of [46] would prove the conjecture. Neither the theorem nor the proposed argument is included, so the conjecture is not marked proved.

The trailing context says the non-Hausdorff quotient \(\operatorname{PF}(F)/\operatorname{PMC}(F)\) contains \(\operatorname{Arc}(F)\) as an open dense stratified subspace, connecting it with a classical, nonquantum aspect of Problem 1. The notation and Problem 1 are absent. It suggests techniques of [1,34] for parts of Problem 6. These are contextual source claims, not additional definitions or verified computations.
\subsection{Short English statement}
Compute the topology of the source's stratified arc complexes, beginning with four specified nonspherical manifold cases of dimensions \(5,7,7,9\), and investigate the conjecture that every arc complex has finite recursively defined type. The four surfaces and precise arc-space construction are omitted.
\subsection{Sources}
\begin{itemize}
\item Local UnsolvedMath registry, permanent ID 11000136, source code catalogue entry, Problem 6, PDF page 118. The chapter header identifies R. Penner; no individual proposer field is supplied.
\item Farb, \emph{Problems on Mapping Class Groups and Related Topics}. Supplied URL: \url{https://www.math.uchicago.edu/~farb/papers/mcgbook.pdf}; recorded access 2026-07-29; extraction pdf-font-structure. The volume contains 329 explicitly typeset headings across contributed chapters, with repeated numbering disambiguated by page.
\item Internal Theorem 5, Problem 1, reference [46] on a mapping-cylinder collapse, and [1,34] for possible techniques are cited without the corresponding statements or full bibliography. The research summary also mentions Penner--Harer and decorated Teichm\"uller theory without precise citations.
\item Supplied status open and OPEN-TRIAGE. The summary reports no computation of the four type-one topological types and no proof or disproof of the finite-type conjecture located in its triage. This is a supplied search assessment. Current primary-literature status and redistribution terms remain NEEDS\_REVIEW.
\end{itemize}
\end{document}
